% Figure 13-A -- CCB-D block diagram.
% INTAKE -> POLICY -> AMPLIFY -> RHYTHM -> MANIFEST, domain diodes D2/D3 in the
% signal path, ^5^P bias rail along the bottom, REFUSE branch through F1, and the
% two feedback loops (fast FLT1 inside ^O; slow FLT2+FB1 terminating at C1).
\begin{tikzpicture}[
  x=1cm, y=1cm,
  stage/.style={casblock, text width=14.5mm, minimum height=14mm,
                font=\sffamily\scriptsize},
  smalllbl/.style={font=\sffamily\scriptsize, align=center, inner sep=1pt},
  ilbl/.style={font=\sffamily\scriptsize\itshape, align=center, inner sep=1pt},
  loopbox/.style={nblock, font=\sffamily\scriptsize, minimum height=6mm,
                  inner sep=2.5pt},
  % Domain diode: triangle + bar, current left to right.
  diodesym/.pic={
    \filldraw[fill=white, draw=black, line width=0.6pt]
      (-0.19,-0.18) -- (-0.19,0.18) -- (0.12,0) -- cycle;
    \draw[line width=0.9pt] (0.12,-0.21) -- (0.12,0.21);
  },
  % Fuse: rectangle with the conductor through it (vertical orientation).
  fusesym/.pic={
    \draw[line width=0.6pt] (0,-0.24) -- (0,0.24);
    \filldraw[fill=white, draw=black, line width=0.6pt]
      (-0.13,-0.17) rectangle (0.13,0.17);
    \draw[line width=0.6pt] (0,-0.17) -- (0,0.17);
  }
]

% ---------------- main signal path ----------------
\node[stage, cblock] (intake)   at (0.83,0) {\textbf{INTAKE}\\
      \casnote{\detokenize{^0^C}}\\ R1\\ k .30--.64};
\node[stage, pblock] (policy)   at (2.81,0) {\textbf{POLICY}\\
      \casnote{\detokenize{^3^P}}\\ Q0 $\cdot$ Q1\\ sat 0.85};
\node[stage, eblock] (amplify)  at (5.09,0) {\textbf{AMPLIFY}\\
      \casnote{\detokenize{^9^E}}\\ Q2\\ sat 0.80};
\node[stage, eblock] (rhythm)   at (7.07,0) {\textbf{RHYTHM}\\
      \casnote{\detokenize{^7^E}}\\ OSC1\\ 24\,hr};
\node[stage, oblock, text width=17mm, font=\sffamily\scriptsize]
      (manifest) at (9.35,0) {\textbf{MANIFEST}\\
      \casnote{\detokenize{^9^O}}\\ Q3\\ sat 1.0};

\draw[flow] (1.66,0) -- (1.98,0);
\draw[flow] (5.92,0) -- (6.24,0);

% domain diodes in the path
\draw[line width=0.6pt] (3.64,0) -- (4.26,0);
\pic at (3.95,0) {diodesym};
\draw[flow] (4.12,0) -- (4.26,0);
\node[smalllbl, anchor=south] at (3.95,0.85)
      {\textbf{D2}\\ \casnote{\detokenize{^C -> ^E}}\\ Vf 1.2};

\draw[line width=0.6pt] (7.90,0) -- (8.52,0);
\pic at (8.21,0) {diodesym};
\draw[flow] (8.38,0) -- (8.52,0);
\node[smalllbl, anchor=south] at (8.21,0.85)
      {\textbf{D3}\\ \casnote{\detokenize{^E -> ^O}}\\ Vf 2.0};

% input / output terminals
\node[casterm, anchor=west] at (0,2.05) {\casnote{\detokenize{^0_desire^C_{F2}|A}}};
\draw[flow] (0.83,1.75) -- (0.83,0.70);
\draw[flow] (10.18,0) -- (10.40,0);
\node[smalllbl, anchor=east] at (10.40,1.05) {TP4\\ \casnote{\detokenize{^9^O}}};

% ---------------- REFUSE branch through F1 ----------------
\draw[flow, rounded corners=2pt] (3.36,-0.70) -- (3.36,-1.15) -- (4.22,-1.15)
      -- (4.22,-1.93);
\pic at (4.22,-1.50) {fusesym};
\node[smalllbl, anchor=west] at (4.40,-1.50) {\textbf{F1} 0.90};
\node[pblock, text width=17mm, font=\sffamily\scriptsize] (refuse) at (4.22,-2.28)
      {\textbf{REFUSE}\\ permanent block};

% ---------------- bias rail + C1 ----------------
\draw[line width=1.3pt] (0.55,-3.15) -- (5.50,-3.15);
\foreach \x/\n in {2.15/Q0, 2.90/Q1, 5.50/Q2}{
  \draw[line width=0.6pt] (\x,-3.15) -- (\x,-0.70);
  \fill (\x,-3.15) circle (0.05);
  \node[smalllbl, anchor=south west] at (\x+0.06,-3.10) {\n};
}
\node[smalllbl, anchor=west] at (6.10,-3.15)
      {\casnote{\detokenize{^5^P}} bias rail};

% C1 hanging off the left end of the rail
\draw[line width=0.6pt] (0.55,-3.15) -- (0.55,-3.45);
\draw[line width=0.9pt] (0.32,-3.45) -- (0.78,-3.45);
\draw[line width=0.9pt] (0.32,-3.62) -- (0.78,-3.62);
\draw[line width=0.6pt] (0.55,-3.62) -- (0.55,-3.90);
\draw[line width=0.7pt] (0.35,-3.90) -- (0.75,-3.90);
\draw[line width=0.7pt] (0.42,-4.00) -- (0.68,-4.00);
\draw[line width=0.7pt] (0.49,-4.10) -- (0.61,-4.10);
\node[smalllbl, anchor=west, align=left] at (0.92,-3.62)
      {\textbf{C1}\ \casnote{\detokenize{^5_operator^P}}\\ $\tau$=3333\,hr};

% ---------------- LOOP A: fast rate limiter, stays inside ^O ----------------
\draw[fbflow, rounded corners=2pt] (9.35,-0.70) -- (9.35,-1.85) -- (7.07,-1.85)
      -- (7.07,-0.70);
\node[loopbox] at (8.21,-1.85) {FLT1\quad $\tau$=24\,hr};
\node[ilbl, anchor=north] at (8.21,-2.18)
      {re-trigger limit, stays in \casnote{\detokenize{^O}}};

% ---------------- LOOP B: slow, lawful ^O -> ^P path ----------------
\fill (10.30,0) circle (0.05);
\draw[fbflow, rounded corners=2pt] (10.30,0) -- (10.30,-4.80) -- (8.75,-4.80);
\draw[fbflow] (6.45,-4.80) -- (5.20,-4.80);
\draw[fbflow, rounded corners=2pt] (3.40,-4.80) -- (0.15,-4.80) -- (0.15,-3.15)
      -- (0.50,-3.15);
\node[loopbox] at (7.60,-4.80) {FLT2\quad $\tau$=3\,weeks};
\node[loopbox] at (4.30,-4.80) {FB1\quad $\beta$=0.05};
\node[ilbl, anchor=south] at (7.60,-4.48)
      {lawful \casnote{\detokenize{^O}}$\,\to\,$\casnote{\detokenize{^P}}};
\node[ilbl, anchor=north] at (4.30,-5.12)
      {terminates at the bias terminal, never at the signal path};

\end{tikzpicture}
